%% kapa: exam=1 round=22 year=2011 period=1 form=A subject=civil_law subject_ko=민법 seq=35 number=35
%% source: Q-Net 감정평가사 기출문제 (cst00309) artlSeq=5036790 file=extracted/1008/5036790/2011년 제22회 감정평가사 1차시험 기출문제/1교시.hwp
%% conversion: kordoc 4.24.51 -> tools/gen1.py (auto); review=proofread-sonnet; flags=none
\begin{question}{35}
합유에 관한 설명으로 옳지 않은 것은? (다툼이 있으면 판례에 의함)
\begin{choices}
  \item 합유자는 합유물의 분할을 청구하지 못한다.
  \item 합유자는 전원의 동의 없이 합유물에 대한 지분을 처분하지 못한다.
  \item 합유는 조합체의 해산 또는 합유물의 양도로 인하여 종료한다.
  \item 합유자는 합유물에 관하여 경료된 원인무효의 소유권이전등기의 말소를 단독으로 청구할 수 없다.
  \item 부동산의 합유자 중 일부가 사망한 경우, 합유자 사이에 특별한 약정이 없는 한 사망한 합유자의 상속인은 합유자로서의 지위를 승계하지 못한다.
\end{choices}
\end{question}
\answer{4}
